\section{Retina and iris profiles}
\label{sec:s-profiles}
The brain and face profiles and the maps of every dataset are in the main paper.
Tables~\ref{tab:fundus} and~\ref{tab:iris} give the retina and iris on the scale $Q$ of Eq.~(1) of the main paper (1 = dataset mean, about 2 = stranger); bold marks the fill closest to the target and underline the second closest; lower is closer.

\begin{table}[htbp]
\centering
\caption{\textbf{Retina, provisional.}
Context mean is the mean of 16 look-alikes; intervals resample people; PAPILA has no second visit.
The fellow eye is further than the dataset mean in both datasets; the second visit is clearly closest, and the context mean is slightly closer than the dataset mean.
Source: \src{results/donor_ladder_2d_facefmt.json}.}
\label{tab:fundus}
\small
\begin{tabular}{@{}lcccc@{}}
\toprule
Dataset & Second visit & Fellow eye & Context mean & Stranger \\
\midrule
GRAPE P8 & \shortstack{\textbf{0.489}\statuspre\\{\scriptsize [0.404, 0.585]}} & \shortstack{1.330\statuspre\\{\scriptsize [1.224, 1.437]}} & \shortstack{\underline{0.856}\statuspre\\{\scriptsize [0.832, 0.880]}} & \shortstack{2.012\statuspre\\{\scriptsize [1.897, 2.140]}} \\
PAPILA P9 & \textit{unavailable} & \shortstack{\underline{1.467}\statuspre\\{\scriptsize [1.422, 1.512]}} & \shortstack{\textbf{0.935}\statuspre\\{\scriptsize [0.914, 0.955]}} & \shortstack{1.992\statuspre\\{\scriptsize [1.941, 2.049]}} \\
\bottomrule

\end{tabular}
\end{table}

\begin{table}[htbp]
\centering
\caption{\textbf{Iris, provisional.}
Look-alike is the single unrelated eye with the closest context; intervals resample people.
Target eyes and people: 156 and 78 (UTIRIS \acs{NIR}), 158 and 79 (UTIRIS \acs{VIS}), 128 and 64 (UPOL).
In all three sets the same eye captured again is closest, a look-alike comes second, and the person's other eye is further than the dataset mean.
Source: \src{results/iris_ladder_metrics.json}.}
\label{tab:iris}
\small
\begin{tabular}{@{}lcccc@{}}
\toprule
Dataset & Same eye again & Other eye & Look-alike & Stranger \\
\midrule
UTIRIS NIR & \shortstack{\textbf{0.407}\statuspre\\{\scriptsize [0.345, 0.475]}} & \shortstack{1.256\statuspre\\{\scriptsize [1.123, 1.395]}} & \shortstack{\underline{0.599}\statuspre\\{\scriptsize [0.540, 0.660]}} & \shortstack{2.001\statuspre\\{\scriptsize [1.816, 2.207]}} \\
UTIRIS VIS & \shortstack{\textbf{0.569}\statuspre\\{\scriptsize [0.503, 0.640]}} & \shortstack{1.293\statuspre\\{\scriptsize [1.211, 1.375]}} & \shortstack{\underline{0.861}\statuspre\\{\scriptsize [0.826, 0.898]}} & \shortstack{2.018\statuspre\\{\scriptsize [1.884, 2.176]}} \\
UPOL & \shortstack{\textbf{0.177}\statuspre\\{\scriptsize [0.154, 0.200]}} & \shortstack{1.118\statuspre\\{\scriptsize [1.037, 1.198]}} & \shortstack{\underline{0.841}\statuspre\\{\scriptsize [0.798, 0.882]}} & \shortstack{2.138\statuspre\\{\scriptsize [2.007, 2.281]}} \\
\bottomrule

\end{tabular}
\end{table}
